\chapter{Diophantine Equations}\label{ch:diophantine}

\begin{center}\itshape
$3^2+4^2=5^2$, but $x^4+y^4=z^4$ has no solution, and neither does $y^2=x^3+7$.\\
Four ideas solve a surprising number of equations: factor, reduce modulo $n$, draw a line, descend.
\end{center}


The equation $x^2+y^2=5$ has the integer solutions $(\pm1,\pm2)$ and $(\pm2,\pm1)$. The equation $x^2+y^2=3$ has no integer solution. To prove the second claim, reduce modulo $4$: an integer square is $0$ or $1$, and two such squares cannot add to $3$. One equation asks us to describe a family of solutions; the other asks us to prove that a family is empty.

This is the first choice in a Diophantine problem: what are the allowed numbers? A real solution of $x^2+y^2=3$ exists, but that does not settle the integer question. If the variables must be positive, the signs matter. If they may be zero, a factorisation may have extra solutions. State the domain before the clever part.

Four recurring methods will do most of the work here. A congruence can forbid a solution. A factorisation can reduce infinitely many candidates to finitely many divisors. A line through a known rational point can describe further rational points on a conic. Infinite descent can turn a supposed smallest solution into an even smaller one, which is impossible.

These methods will not solve every integer equation. The absence of a universal algorithm is a deep theorem, not an excuse to stop thinking. We will choose methods because of the structure of the equation, check their hypotheses, and say clearly what they prove.


\subsection*{Choose the battlefield before choosing the weapon}
An equation comes with a domain. In non-negative integers, $x^4+y^4=z^2$ has the uninteresting solutions with one of $x,y$ zero; in \emph{positive} integers Fermat's theorem excludes all solutions. In rational numbers $x^2+y^2=3$ is a different question from its solvability in real numbers. Write down the domain first: it determines which divisions, order arguments and congruences are available.

Here is a practical sequence of questions. What happens modulo $4$, $8$ or a prime dividing a coefficient? Can we divide out a common factor? Can a product be split into coprime factors? Is the equation quadratic in one variable, allowing a second root? Is there a natural positive integer that should decrease? These are prompts, not an algorithm guaranteed to finish.

A search has a different job. Finding a solution certifies existence; finding none up to a million certifies only that particular finite range. Number theory keeps a large drawer for plausible statements which survived an impressive search and were still false. Proof is the drawer's lock.

\section{Obstructions before constructions}

% t1-pass2 motive 8dd733c8
Before searching for a solution one should ask whether a cheap congruence already forbids one. A single prime modulus can finish the problem, and the same idea, completed at every prime and at the real place, becomes the Hasse principle of Volume~13. The point of this section is to know when an obstruction is a proof and when it is only a hint.

\subsection*{A congruence can finish the problem}
The equation $y^2=3x^2+2$ has no integer solution: reducing modulo $3$
would make a square equal to $2$, whereas the only squares modulo $3$ are
$0,1$. The negative conclusion is complete. It applies to every integer
$x$, including negative ones, because reduction forgets the sign.

A successful congruence test has a different meaning. For $y^2=3x^2+1$,
modulo $3$ merely requires $y\equiv\pm1$. It leaves infinitely many
candidates, most of which need not solve the equation. For example $(x,y)=(1,1)$
passes this test but fails the original equation. Nevertheless $(1,2)$ is a
solution. Congruences are necessary conditions; making them sufficient needs
an additional theorem or a construction.

\subsection*{Valuations detect an imbalance}
Suppose $x^2=12y^2$. If $y=0$, then $x=0$. Otherwise neither side is zero,
and taking the $3$-adic valuation gives
\[
2v_3(x)=1+2v_3(y),
\]
an even integer equal to an odd one. Thus $(0,0)$ is the only integer
solution. Taking the $2$-adic valuation instead gives
$2v_2(x)=2+2v_2(y)$, which has no parity contradiction. The prime must be
chosen to reveal the obstruction.

More generally, let $k\ge1$ be an integer. If $A,B\ne0$ and $Ax^k=By^k$ with $xy\ne0$, then
\[
v_p(A)-v_p(B)=k\bigl(v_p(y)-v_p(x)\bigr)
\]
for every prime $p$. A difference not divisible by $k$ rules out a solution.
The same test works over $\Q$, where valuations may be negative.
It says nothing about a sum until cancellation has been understood.
For $x^2=9y^2$ all these exponent conditions pass, and factorization gives
$x=3y$ or $x=-3y$. Passing the test is a reason to change tools.

\subsection*{Factorization makes a finite search complete}
Find the solutions of $x^2-y^2=45$ with $x>y\ge0$. Put
$u=x-y$, $v=x+y$. Then $uv=45$, $0<u\le v$, and $u,v$ have the same parity.
Conversely any positive factor pair with that parity gives
\[
x=(u+v)/2,\qquad y=(v-u)/2.
\]
The three pairs $(1,45),(3,15),(5,9)$ yield
$(x,y)=(23,22),(9,6),(7,2)$, and these are all the solutions.
This search is a proof because factorization bounds and exhausts its candidates.
An unexplained search over $x\le1000$ would have a different logical status.

The factors need not be coprime: the pair $(3,15)$ has gcd $3$.
A perfect-power equation requires extra care before splitting its factors
into powers. The Pythagorean parametrization below supplies the model:
remove a common factor, check parity, and only then use prime exponents.

\begin{think}
For $x^2-y^2=12$, must the two factors be odd or even?
The pair $(2,6)$ gives $(4,2)$; the pairs $(1,12)$ and $(3,4)$ fail the
parity condition. What changes for $x^2-y^2=2$?
\end{think}

% t1-pass2 examples 8dd733c8
\begin{examplex}[no solution modulo $8$]
The equation $y^2=x^2+x+1$ looks harmless, but modulo $2$ it says a square is $1$ when $x$ is even and $1$ when $x$ is odd, so that test passes. Modulo $4$, squares are $0$ or $1$. If $x$ is even, $x=2t$, then $y^2\equiv 0+0+1=1$. If $x$ is odd, $x^2+x\equiv 1+1\equiv 2\pmod 4$, so $y^2\equiv 3\pmod 4$, which is impossible. Odd $x$ is excluded; even $x$ is not, and indeed $(x,y)=(0,1)$ works. The modulus eliminated half the problem and left a real solution standing.
\end{examplex}

\begin{examplex}[a modulus that finishes the job]
Consider $y^2=3x^2+2$. Modulo $3$ this says $y^2\equiv 2$. The squares modulo $3$ are $0$ and $1$, so there is no integer solution at all. The same obstruction, rewritten as $y^2-3x^2=2$, is the sort of local condition Volume~13 organises for every conic.
\end{examplex}

\section{Pythagorean triples}

% t1-pass2 motive 049c5671
The equation $x^2+y^2=z^2$ is the first Diophantine equation with infinitely many solutions and a complete list. The list is a multiplication in $\Z[i]$ written without Gaussian integers, and it is the input to Fermat's descent for fourth powers. Volume~13's congruent-number problem asks which of these triples have square area; that question is elliptic, and it begins here.

A solution of $x^2+y^2=z^2$ in positive integers is \emph{primitive} if $\gcd(x,y,z)=1$. Every solution is a multiple of a primitive one (divide by the gcd), and in a primitive one $x,y$ are coprime (a prime dividing both divides $z^2$, hence $z$).

\begin{think}
The primitive triples with $z<30$ are $(3,4,5)$, $(5,12,13)$, $(15,8,17)$, $(7,24,25)$ and $(21,20,29)$, written with $x$ odd. For each, compute $\frac{z+x}2$ and $\frac{z-x}2$. What do you notice? Can you explain it from $y^2=(z-x)(z+x)$?
\end{think}

\begin{theorem}\label{thm:pyth}
The primitive solutions of $x^2+y^2=z^2$ with $y$ even are exactly
\[
x=m^2-n^2,\qquad y=2mn,\qquad z=m^2+n^2,
\]
with $m>n\ge1$ coprime and of opposite parity.
\end{theorem}
\begin{proof}
These are solutions: $(m^2-n^2)^2+(2mn)^2=(m^2+n^2)^2$. They are primitive: a prime $p$ dividing $x$ and $z$ divides $2m^2$ and $2n^2$; $p\ne2$ since $x$ is odd ($m,n$ of opposite parity), so $p\mid m$ and $p\mid n$, impossible.

Conversely, let $(x,y,z)$ be primitive. $x,y$ are not both odd (else $z^2\equiv2\pmod4$), and not both even, so exactly one is even, say $y$; then $x,z$ are odd. Write $y^2=(z-x)(z+x)$ and $\left(\frac y2\right)^2=\frac{z-x}2\cdot\frac{z+x}2$, a product of two integers. They are coprime: a common divisor divides their sum $z$ and difference $x$. A product of two coprime positive integers is a square only if each is a square (compare the valuations), so $\frac{z+x}2=m^2$ and $\frac{z-x}2=n^2$ with $m>n\ge1$, $\gcd(m,n)=1$. Then $z=m^2+n^2$, $x=m^2-n^2$, $y=2mn$, and $m,n$ have opposite parity because $z$ is odd.
\end{proof}

% t1-pass2 why
\paragraph{Why this works.}
A primitive triple has exactly one even leg. The two factors $(z-x)/2$ and $(z+x)/2$ are then coprime positive integers whose product is a square, so each is a square, say $n^2$ and $m^2$. Opposite parity of $m$ and $n$ is the condition that $z$ is odd. Every primitive triple arises once this way.

The key step, ``coprime factors of a square are squares'', is the first example of the \emph{factorisation method}. It will be used again with more exotic factorisations in Chapter~\ref{ch:gaussian}.

\subsection*{Using the parameters, and checking their restrictions}
Take $(m,n)=(4,1)$. The formula yields $(15,8,17)$, and the inverse calculation recovers $m^2=(17+15)/2=16$ and $n^2=(17-15)/2=1$. Thus the parameters are not lucky guesses: the factorisation forces them.

If $m,n$ are both odd, the three displayed coordinates are all even: $(3,1)$ gives $(8,6,10)$, a doubled $(4,3,5)$. If they share a factor $d$, all three coordinates share $d^2$. This explains both restrictions in the theorem. Conversely, they guarantee primitivity, so we can generate every integer right triangle by taking a permitted pair and then multiplying its three sides by an arbitrary positive integer.

As a small geometric payoff, the inradius $r$ of such a primitive right triangle is
\[
r=\frac{x+y-z}{2}=n(m-n).
\]
Indeed, its area is $xy/2$ and its semiperimeter is $(x+y+z)/2$; their quotient is $xy/(x+y+z)=(x+y-z)/2$ because $x^2+y^2=z^2$. With an overall scale factor $k$, the inradius is $kn(m-n)$. This provides a useful independent check on generated triples.

\begin{figure}[ht]
\centering
\fig{ch09-triangle-incircle}
\caption{The $3$--$4$--$5$ triangle has area $6$, semiperimeter $6$, and inradius $1$. The three dashed radii meet the sides at right angles. The formula $r=(a+b-c)/2$ connects the geometry with the integer parameters of a Pythagorean triple.}
\label{fig:incircle}
\end{figure}

% t1-pass2 examples 049c5671
\begin{examplex}[two triples from the formulas]
For $m=2$, $n=1$, one gets $x=3$, $y=4$, $z=5$.
For $m=4$, $n=1$, one gets $x=15$, $y=8$, $z=17$, and $15^2+8^2=225+64=289=17^2$.
Both pairs $(m,n)$ are coprime and of opposite parity, so both triples are primitive.
\end{examplex}

\begin{examplex}[a pair the formulas reject]
The choice $m=3$, $n=1$ is legal: $x=8$, $y=6$, $z=10$, which is $2\cdot(3,4,5)$, not primitive, because $m$ and $n$ are both odd. Opposite parity was the condition that kept $x$ odd and the triple primitive. The formulas do not miss the imprimitive triples; those are integer multiples of primitive ones.
\end{examplex}

\section{Infinite descent}

% t1-pass2 motive b63dce63
Fermat's proof that $x^4+y^4=z^2$ has no positive solution is descent: a solution would manufacture a smaller positive solution. There is no infinite decreasing sequence of positive integers, so there is no solution. Figure~\ref{fig:t1-descent} is the shape of the argument. The same shape, rewritten with the group law, is Mordell's theorem in Volume~13.

\begin{theorem}[Fermat]\label{thm:fermat4}
The equation $x^4+y^4=z^2$ has no solution in positive integers. In particular $x^4+y^4=z^4$ has none.
\end{theorem}
\begin{proof}
Suppose there are solutions, and take one with $z$ minimal. Then $\gcd(x,y)=1$: if a prime $p$ divides $x$ and $y$, then $p^4\mid z^2$, so $p^2\mid z$, and $\left(\frac xp,\frac yp,\frac z{p^2}\right)$ is a smaller solution. So $(x^2,y^2,z)$ is a primitive Pythagorean triple; say $y$ is even. By Theorem~\ref{thm:pyth},
\[
x^2=m^2-n^2,\qquad y^2=2mn,\qquad z=m^2+n^2,
\]
with $m>n\ge1$ coprime of opposite parity. From $x^2+n^2=m^2$ with $x$ odd, $(x,n,m)$ is a primitive Pythagorean triple (any common prime of $n,m$ would divide $\gcd(m,n)=1$), so $n$ is even (if $n$ were odd, $m^2=x^2+n^2\equiv2\pmod4$), and $m$ is odd, and $n=2ab$, $m=a^2+b^2$ with $a,b$ coprime. Now $y^2=2mn=4mab$, so $\left(\frac y2\right)^2=mab$, where $m,a,b$ are pairwise coprime ($a,b$ divide $n$, which is prime to $m$). So each of $m,a,b$ is a square: $a=u^2$, $b=v^2$, $m=w^2$. Then
\[
u^4+v^4=a^2+b^2=m=w^2,
\]
a new solution with $w\le w^2=m\le m^2<m^2+n^2=z$. This contradicts the minimality of $z$.
\end{proof}

% t1-pass2 why
\paragraph{Why this works.}
A solution of $x^4+y^4=z^2$ is a Pythagorean triple whose legs are squares. The parameters $m>n$ of that triple are forced to be squares themselves, and they build a new positive solution whose hypotenuse is strictly smaller. No infinite descent of positive integers exists, so the first solution did not exist either.

\begin{remark}
Fermat's Last Theorem, that $x^n+y^n=z^n$ has no solution in positive integers for $n\ge3$, was proved by Wiles (with Taylor) in 1995. Theorem~\ref{thm:fermat4} settles every $n$ divisible by $4$; Euler's case $n=3$ needs the Eisenstein integers of Chapter~\ref{ch:gaussian}. The general proof is far beyond this series, and we never use it.
\end{remark}

\subsection*{What makes descent a proof?}
A descent requires an assumed solution, a rule producing another solution, and a \emph{strictly smaller positive integer}. Here the monitored integer is $z$. We do not assume that repeatedly drawing smaller triangles somehow disappears; we use the well-ordering principle to choose a least $z$ and contradict its minimality in one step.

The stronger equation $x^4+y^4=z^2$ is chosen deliberately. Descending while insisting on $z^4$ would impose a condition the construction need not preserve. By ruling out the square on the right, we automatically rule out a fourth power, since $z^4=(z^2)^2$. Strengthening a claim can sometimes make its proof easier: the stronger form is stable under the operation we need.

The sentence ``coprime factors of a square are squares'' is also doing real work. The product $2\cdot8$ is a square although neither factor is. In the proof, $m,a,b$ are pairwise coprime, so every prime exponent in the product belongs to a single factor and must be even. Without that check, the descent would have a hole large enough for Fermat's missing margin.

% t1-pass2 examples b63dce63
\begin{examplex}[the first step on a hypothetical solution]
Suppose $x^4+y^4=z^2$ with $x,y,z>0$. Then $(x^2,y^2,z)$ is a Pythagorean triple. If it is primitive and $y$ is even, the parametrization gives $x^2=m^2-n^2$ and $y^2=2mn$ with $m>n$. The descent shows that $m$ and $n$ are themselves squares of a smaller solution. One never needs a numerical example of a solution, because the argument shows there is nothing to plug in.
\end{examplex}

\begin{figure}[ht]
\centering
\fig{ch09-descent}
\caption{The shape of Fermat's descent. A positive solution of $x^4+y^4=z^2$ is a Pythagorean triple with two square legs, and the standard parameters of that triple are a strictly smaller positive solution.}
\label{fig:t1-descent}
\end{figure}

\begin{examplex}[descent on a real equation, small enough to finish]
The equation $x^4+y^4=z^4$ is the case $z^2$ replaced by a fourth power, so it is included. For a contrast that does have solutions, $x^2+y^2=z^2$ has no such descent: the parametrization produces a finite pair $(m,n)$, not another triple of the same shape. Descent requires the new solution to satisfy the original equation.
\end{examplex}

\section{Rational points on conics: the chord method}

% t1-pass2 motive c6dd1f1c
One rational point on a conic produces all the others, by drawing lines of rational slope. The algebra is that a quadratic polynomial with one rational root has another. This is the geometric form of the parametrization of Pythagorean triples, and it fails for cubics, where a line meets the curve three times and the set of rational points becomes a group. That group is the subject of Volume~13.

\begin{think}
Draw the line through $(-1,0)$ with slope $\frac12$. Where does it meet the circle $x^2+y^2=1$ again? Before computing: why must that second point have rational coordinates? (Think about the sum of the two roots of a quadratic with rational coefficients.)
\end{think}

\begin{proposition}\label{prop:chord}
Let $Q\in\Q[x,y]$ have degree $2$ and no factor of degree $1$ in $\Q[x,y]$, and suppose the conic $Q(x,y)=0$ contains a rational point $P=(x_0,y_0)$. Then every line through $P$ with rational slope $t$ meets the conic in at most one other point, which is rational. Conversely, every rational point other than $P$ is obtained this way, from the line through it and $P$ (or from the vertical line through $P$).
\end{proposition}
\begin{proof}
Substitute $x=x_0+s$, $y=y_0+ts$ in $Q$: we get $g(s)=\alpha s^2+\beta s+\gamma$ with $\alpha,\beta,\gamma\in\Q$, and $\gamma=Q(x_0,y_0)=0$. The polynomial $g$ is not zero: otherwise $Q$ vanishes on the whole line, and writing $\ell=y-y_0-t(x-x_0)$ and $Q=\ell A+Q\big(x,y_0+t(x-x_0)\big)$ (divide by $\ell$ as a polynomial in $y$), the remainder vanishes identically, so $\ell\mid Q$, which we excluded. (Without this hypothesis the statement fails: on $xy=0$ the line $y=0$ through $(0,0)$ lies entirely in the conic.) So $g(s)=s(\alpha s+\beta)$ with $(\alpha,\beta)\ne(0,0)$, whose roots are $0$ (the point $P$) and, if $\alpha\ne0$, $s=-\frac\beta\alpha\in\Q$, giving a rational point. Conversely, if $P'\ne P$ is rational, the line $PP'$ has rational slope (or is vertical), and $P'$ is its second intersection.
\end{proof}

% t1-pass2 why
\paragraph{Why this works.}
Substitution of a rational line into a quadratic produces a quadratic polynomial in one variable. One root is known and rational, the coefficients are rational, and the product of the roots is rational, so the second root is rational. That second intersection is the new point. A vertical line is the same argument with the variables swapped.

\begin{figure}[ht]
\centering
\fig{ch09-chord-markov}
\caption{Left: the chord method on the unit circle. Each line through $(-1,0)$ with rational slope $t$ meets the circle again at a rational point, and every rational point arises this way; $t=\frac12$ gives the triple $(3,4,5)$, $t=\frac23$ gives $(5,12,13)$. Right: the first levels of the tree of Markov triples (Theorem~\ref{thm:markov}). Each triple $(a,b,c)$ with $c$ largest has two children, $(a,c,3ac-b)$ and $(b,c,3bc-a)$.}
\label{fig:chordmarkov}
\end{figure}

\begin{example}
On $x^2+y^2=1$ with $P=(-1,0)$, the line $y=t(x+1)$ gives $x^2-1+t^2(x+1)^2=0$, i.e.\ $(x+1)\left((x-1)+t^2(x+1)\right)=0$. So the rational points are $(-1,0)$ and
\[
\left(\frac{1-t^2}{1+t^2},\ \frac{2t}{1+t^2}\right),\qquad t\in\Q.
\]
With $t=\frac nm$ and clearing denominators, this recovers Theorem~\ref{thm:pyth}. It is also the substitution $t=\tan\frac\theta2$ of calculus.
\end{example}

A conic without rational points can often be detected by a congruence: $x^2+y^2=3$ has no rational point, because $X^2+Y^2=3Z^2$ has no non-zero integer solution (Exercise~\ref{ex:x2y23z2}). Legendre's theorem (a special case of the Hasse--Minkowski theorem, named at the end of the $p$-adic chapter of Volume~13 and proved there) says that congruence obstructions are the only ones for conics; for cubics this fails (Selmer's $3x^3+4y^3+5z^3=0$).

% t1-pass2 examples c6dd1f1c
\begin{examplex}[the circle through $(1,0)$]
On $x^2+y^2=1$, the point $P=(1,0)$ is rational. The line of slope $t$ through $P$ is $y=t(x-1)$. Substituting and discarding the known root $x=1$ leaves the second intersection
\[
x=\frac{t^2-1}{t^2+1},\qquad y=\frac{-2t}{t^2+1}
\]
after a short calculation, up to the sign convention for the slope. For $t=2$ one gets a rational point, and clearing denominators recovers a Pythagorean triple.
\end{examplex}

\begin{examplex}[a conic with no rational point]
The conic $x^2+y^2=3$ has real points and no rational point. If $X^2+Y^2=3Z^2$ had a nonzero integer solution, reducing modulo $4$ would be impossible for a primitive solution, since squares are $0$ or $1$ modulo $4$ and the left side cannot be $3$ times a square. The chord method never starts, because it has no rational point to start from.
\end{examplex}

\section{Factoring with a sum of squares}

% t1-pass2 motive a0518b8e
Some cubic equations become a sum of two squares after a shift, and a prime $3$ modulo $4$ is then not allowed to odd valuation. The method is short and local. It is enough for $y^2=x^3+7$, and it is the ancestor of the descent on elliptic curves.

\begin{theorem}[a Mordell equation]\label{thm:mordell7}
The equation $y^2=x^3+7$ has no solution in integers.
\end{theorem}
\begin{proof}
If $x$ is even, $y^2\equiv7\pmod8$, impossible since squares are $0,1,4$ modulo $8$. So $x$ is odd. Add $1$ to both sides:
\[
\begin{gathered}y^2+1=x^3+8=(x+2)(x^2-2x+4),\qq x^2-2x+4=(x-1)^2+3.\end{gathered}
\]
As $x-1$ is even, $(x-1)^2+3\equiv3\pmod4$, and it is $\ge3$. A positive integer $\equiv3\pmod4$ has a prime factor $q\equiv3\pmod4$ (a product of primes all $\equiv1\pmod4$ is $\equiv1$, and it is odd). Then $q\mid y^2+1$, so $-1$ is a square modulo $q$, which contradicts $q\equiv3\pmod4$ (Theorem~\ref{thm:eulercrit}).
\end{proof}

% t1-pass2 why
\paragraph{Why this works.}
After the shift, $y^2+1$ factors, and one factor is forced to be divisible by a prime $q\equiv 3\pmod 4$ to an odd power. But $y^2+1$ is a sum of two squares, so Theorem~\ref{thm:twosqn} forbids exactly that. The congruence did not merely fail for small $x$; it fails for every integer $x$.

The same idea, ``add something to make a sum of two squares, and find a factor $\equiv3\pmod4$'', disposes of many \emph{Mordell equations} $y^2=x^3+k$ (Exercise~\ref{ex:mordell5}). Others, such as $y^2=x^3-2$ (whose only solutions are $(3,\pm5)$), need the arithmetic of $\Z[\sqrt{-2}]$ (Chapter~\ref{ch:gaussian}).

% t1-pass2 examples a0518b8e
\begin{examplex}[why $y^2=x^3+7$ fails for small $x$]
For $x=1$, $y^2=8$, not a square. For $x=2$, $y^2=15$. For $x=-1$, $y^2=6$. For $x=-2$, $y^2=-1$. These are experiments. The proof is that an odd $x$ makes $y^2+1$ divisible by a prime $3$ modulo $4$ to an odd power, which is impossible.
\end{examplex}

\begin{examplex}[the shift]
If $x$ is odd, $x=2k+1$, then $x+2=2k+3$ and $x^2-2x+4=(x-1)^2+3\equiv 3\pmod 4$. A positive integer congruent to $3$ modulo $4$ has a prime factor congruent to $3$ modulo $4$. That prime divides $y^2+1$, and the rest of the proof checks that it divides it only once too often.
\end{examplex}

\section{Vieta jumping}

% t1-pass2 motive 60e71bc1
When one variable appears quadratically, the second root given by Vieta is another integer solution, often a smaller one. Repeating the jump descends to a solution small enough to list. The Markov equation is the classical example, and the same involution appears whenever a Diophantine equation is quadratic in each variable separately.

\begin{think}
$(1,1,1)$ solves $x^2+y^2+z^2=3xyz$. Fix $y=z=1$: then $x$ is a root of $t^2-3t+2$. What is the other root? Now fix the two smallest entries of your new triple and do the same. How far can you go, and does the process ever get stuck?
\end{think}

When an equation is quadratic in one variable, the other root of that quadratic, given by Vi\`ete's relations, is a new solution. This lets us move from solution to solution, and usually to descend.

\begin{theorem}[Markov, 1879]\label{thm:markov}
Every solution of $x^2+y^2+z^2=3xyz$ in positive integers is obtained from $(1,1,1)$ by a finite sequence of permutations and moves $(x,y,z)\mapsto(3yz-x,y,z)$.
\end{theorem}
\begin{proof}
For fixed $y,z$, $x$ is a root of $g(t)=t^2-3yz\,t+(y^2+z^2)$. The other root $x'=3yz-x=\frac{y^2+z^2}x$ is a positive integer, so the move sends solutions to solutions, and it is an involution.

Let $(x,y,z)$ be a solution with $x\ge y\ge z$, other than $(1,1,1)$. We show $x'<x$, so that the move strictly decreases the sum $x+y+z$. Since $z\ge1$, $g(y)=2y^2+z^2-3y^2z\le z^2-y^2\le0$. If $g(y)<0$, then $y$ lies strictly between the roots $x'$ and $x$ of $g$, so $x'<y\le x$. If $g(y)=0$, then $z=1$ and $y=z$, so $y=z=1$ and $x^2-3x+2=0$: $x=2$ (as $x=1$ is excluded), and $x'=1<x$.

So from any solution, repeatedly sorting and applying the move decreases $x+y+z$, which cannot go on forever: we reach a solution where the argument fails, which can only be $(1,1,1)$. Reversing the steps (moves are involutions) builds the original solution from $(1,1,1)$.
\end{proof}

% t1-pass2 why
\paragraph{Why this works.}
Vieta supplies a second integer root, and sorting the coordinates so that the jumped variable is the largest makes the sum of the three variables decrease. A positive integer cannot decrease forever, so every solution descends to a solution that cannot jump, which one checks is a permutation of $(1,1,1)$. Running the involution backwards builds every other solution.

The \emph{Markov numbers} $1$, $2$, $5$, $13$, $29$, $34$, $89$, $169$, $194$, \dots form a tree. The \emph{uniqueness conjecture} (Frobenius, 1913), that the largest element determines the triple, is still open.

\subsection*{A jump, a check, and a descent}
Starting from $(1,2,5)$, keep $2$ and $5$ fixed and replace the $1$ by $3\cdot2\cdot5-1=29$. The result is $(29,2,5)$, and
\[
29^2+2^2+5^2=841+4+25=870=3\cdot29\cdot2\cdot5.
\]
Jumping back replaces $29$ by $30-29=1$. The move is its own inverse: growing and descending use the very same formula. Which direction it takes depends on the coordinate chosen.

In the proof we sort the entries and replace a largest one. This is not cosmetic: replacing a small coordinate usually \emph{increases} the sum, as the example shows. The sign test $g(y)\le0$ locates the smaller root without solving a quadratic by radicals. Vieta jumping turns an apparently static equation into a landscape with a guaranteed downhill path.

\begin{think}
From $(1,5,13)$, replace the $5$ and then undo the move. You get $(1,13,34)$ because $3\cdot1\cdot13-5=34$. Which entry must be replaced to descend from $(1,13,34)$? Try the two other choices and see why the sorting instruction matters.
\end{think}

% t1-pass2 examples 60e71bc1
\begin{examplex}[the smallest Markov solution]
The equation $x^2+y^2+z^2=3xyz$ has the solution $(1,1,1)$: both sides equal $3$.
Vieta in $x$, for $y=z=1$, says the two roots add to $3yz=3$ and multiply to $y^2+z^2=2$. The roots of $t^2-3t+2=0$ are $1$ and $2$. So $(2,1,1)$ is the next solution, and $4+1+1=6=3\cdot 2\cdot 1\cdot 1$.
\end{examplex}

\begin{examplex}[one more jump]
From $(2,1,1)$, jump in the large variable. The quadratic $t^2-(3\cdot 1\cdot 1)t+(1+1)=0$ has roots $1$ and $2$, so the jump swaps $1$ and $2$ and does not yet grow. Jumping from a larger known solution such as $(5,2,1)$ produces a strictly larger partner, which is why the tree of solutions is infinite once a single non-permutation move appears. The proof in the text organises all of them from $(1,1,1)$.
\end{examplex}

\section{Workshop: build the descent before writing it}

% t1-pass2 motive 16e4ab94
A descent is easier to discover on a numerical sketch than to write in letters. This workshop is the place where the reader builds one step of Fermat's argument before checking it against the formal proof. The habit is the one Volume~1 asks for: predict, then write.

To discover the Markov argument, begin with a small table of triples rather
than with a finished inequality. The equation, viewed as a quadratic in $x$,
forces a second root $x'=3yz-x$. There are now three questions: is $x'$
an integer, is it positive, and can one choose the replaced coordinate so
that something decreases? The first follows from the displayed expression;
the second from $xx'=y^2+z^2$; the third requires sorting the coordinates.
None is supplied merely by saying ``Vieta jumping''.

Work backwards from the desired inequality $x'<x$. Comparing explicit
radicals would obscure the arithmetic. Instead test the quadratic at $y$:
$g(y)\le0$ locates $y$ between its roots. The equality case leads to
$(2,1,1)$, which must be checked separately. This is a typical auxiliary
construction: introduce a polynomial whose sign answers a question about
an unknown solution.

Now critique an attempted proof: ``Every jump makes a smaller triple, so
all triples come from $(1,1,1)$.'' The first assertion is false:
$(1,2,5)$ jumps to $(29,2,5)$ when the smallest coordinate is replaced.
The corrected proof chooses a largest coordinate and monitors the positive
integer $x+y+z$. Reversibility then gives generation as well as non-existence
of a counterexample. Choosing the operation, choosing a measure, and
proving its strict decrease are separate pieces of discovery.

The same checklist explains Fermat's earlier descent. We strengthen the
right side from a fourth power to a square because the construction preserves
that larger family. We minimize the new hypotenuse, not a visually smaller
coordinate whose behaviour has not been proved. When a descent stalls,
first check whether the chosen family or the chosen measure is wrong.

\subsection*{When elementary methods have reached their boundary}
A conic with one rational point can be parametrized by lines, as proved
above. The construction needs that first point and the exclusion of a line
component; without either, its hypotheses have not been met. A cubic cannot
be handled by copying the same argument: after substituting a line, removing
a known root leaves a quadratic, whose other roots need not be rational.
That degree count explains both the success and the limitation of the method.
It points toward elliptic curves without borrowing their theory.

% t1-pass2 examples 16e4ab94
\begin{examplex}[the Pythagorean step in numbers that are not a counter-example]
Take the triple $(3,4,5)$, so $3^2+4^2=5^2$. Here the legs are not both squares. The parametrization has $m=2$, $n=1$, and $m$ is not a square. The descent step that would require $m$ and $n$ to be squares simply does not apply. The workshop's point is to see which hypothesis failed, not to force a descent that is not there.
\end{examplex}

\begin{examplex}[what a smaller solution would have to satisfy]
If a primitive solution of $x^4+y^4=z^2$ existed with $y$ even, the parameters would satisfy $m=a^2$, $n=b^2$ and $m-n=c^2$ or an arrangement of the same shape, giving $a^4-b^4$ equal to a square and eventually a new solution of the original equation with smaller hypotenuse. Writing those three squares down, even without values, is the discovery. The proof then checks the inequalities $z_1<z$.
\end{examplex}

\begin{beauty}
Where this leads. Pythagorean triples and the descent for fourth powers are Theorem~9.1 and Theorem~9.2 of this chapter, and they are the first page of the arithmetic of elliptic curves in Volume~13. The chord method on a conic is the Hasse principle in its elementary form; the chord method on a cubic is the group law. Vieta jumping remains an elementary technique long after the equations have become modular.
\end{beauty}

\subsection*{Proof workshop: make a nonexistence argument complete}

Suppose integers $x,y$ satisfy $x^2+y^2=3$.
Every even integer has square $0$ modulo $4$,
and every odd integer has square $1$ modulo $4$.
Therefore the left side is congruent to one of
$0,1,2$ modulo $4$. The right side is $3$ modulo
$4$. This is a contradiction, and it applies to
\emph{all} integers $x,y$.

Notice the three ingredients: a complete list of
possible square residues, an equality that preserves
remainders, and a value outside the list. Checking
only $|x|,|y|\le10$ would not be the same proof,
because it would leave infinitely many candidates.
For existence, the direction reverses: the pair
$(1,2)$ verifies $1^2+2^2=5$, and that single
witness settles the existential question for $5$.
To list \emph{all} solutions requires a separate
argument, here the bounds $x^2,y^2\le5$.


\section{Exercises}

\begin{exercise}[$\star$]\label{ex:areaper}
Find all right triangles with integer sides whose area equals their perimeter.
\end{exercise}

\begin{exercise}[$\star$]\label{ex:x2y23z2}
Show that $x^2+y^2=3z^2$ has no solution in integers other than $(0,0,0)$. Deduce that the circle $x^2+y^2=3$ has no rational point.
\end{exercise}

\begin{exercise}[$\star$]\label{ex:catalan23}
Find all $x,y\in\N$ with $2^x-3^y=1$.
\end{exercise}

\begin{exercise}[$\star\star$]\label{ex:mordell5}
Show that $y^2=x^3-5$ has no solution in integers.
\end{exercise}

\begin{exercise}[$\star\star$]\label{ex:eisentri}
Find all rational points of $x^2+xy+y^2=1$, and deduce all triangles with integer sides $a,b,c$ and an angle of $120^\circ$ opposite $c$, that is, $c^2=a^2+ab+b^2$.
\end{exercise}

\begin{exercise}[$\star\star$]\label{ex:markovcoprime}
Show that the three numbers of a Markov triple are pairwise coprime, and list the triples with largest element at most $200$.
\end{exercise}

\begin{exercise}[$\star\star\star$]\label{ex:imo88}
\emph{IMO 1988, Problem 6.} Let $a,b$ be positive integers such that $k=\dfrac{a^2+b^2}{ab+1}$ is an integer. Show that $k$ is a perfect square.
\end{exercise}

\begin{exercise}[$\star$]\label{ex:t1-triple}
Using $m>n\ge 1$ coprime and of opposite parity, write the primitive triples for $(m,n)=(2,1)$, $(4,1)$ and $(3,2)$. Check one of the identities by hand.
\end{exercise}

\begin{beauty}
Every method of this chapter grows into a theory. Factorisation needs rings where it is unique, as in the Gaussian integers, or ideals when it is not, which is the repair Volume~13 gives for $6=2\cdot3=(1+\sqrt{-5})(1-\sqrt{-5})$. Congruences become the $p$-adic numbers of Volume~13, and the principle ``a conic has a rational point if and only if it has one over $\R$ and over every $\Q_p$'' is the Hasse--Minkowski theorem. The chord method on cubics gives the group law on elliptic curves, where Fermat's descent becomes Mordell's theorem: the rational points of an elliptic curve form a finitely generated group (Volume~13).
\end{beauty}

\section{Hints}

\begin{hint}{ex:areaper}
Substitute $c=\frac{ab}2-a-b$ into $a^2+b^2=c^2$ and simplify. The result is linear in each of $a$ and $b$, so it can be written as $(a-\alpha)(b-\alpha)=\beta$ for suitable integers.
\end{hint}
\begin{hint}{ex:x2y23z2}
Work modulo $3$: squares are $0$ or $1$. Show that $x$ and $y$ are multiples of $3$, then $z$ is too, and descend.
\end{hint}
\begin{hint}{ex:catalan23}
For $y\ge1$, reduce modulo $3$ to see that $x$ is even. Then $3^y=(2^{x/2}-1)(2^{x/2}+1)$: two powers of $3$ that differ by $2$.
\end{hint}
\begin{hint}{ex:mordell5}
Add $4$: $y^2+4=x^3-1=(x-1)(x^2+x+1)$. Check $x$ modulo $4$ first, then find a factor $\equiv3\pmod4$.
\end{hint}
\begin{hint}{ex:eisentri}
Use the rational point $(-1,0)$ and the lines $y=t(x+1)$. For the triangles, divide $c^2=a^2+ab+b^2$ by $c^2$.
\end{hint}
\begin{hint}{ex:markovcoprime}
The move $x\mapsto3yz-x$ does not change $\gcd(x,y,z)$, and neither do permutations. \emph{If stuck on pairwise coprimality:} a prime dividing two coordinates also divides the square of the third, hence the third; a common factor of two would become a common factor of all three. Then grow the tree of Figure~\ref{fig:chordmarkov} until the largest entries exceed $200$.
\end{hint}
\begin{hint}{ex:imo88}
Fix $k$ and take a solution $(a,b)$ with $a\ge b$ and $a+b$ minimal. The other root $a'=kb-a=\frac{b^2-k}a$ of $t^2-kbt+b^2-k=0$ is an integer. Show that $a'\ge0$ if $k$ is not a square, and that $a'<b$; what happens when $a'=0$? \emph{A second hint:} if $a'\le-1$, substitute into $a'^2-kba'+b^2-k$ and use $b\ge1$ to show it is positive, contradicting the root equation.
\end{hint}

\begin{hint}{ex:t1-triple}
The formulas are $m^2-n^2$, $2mn$ and $m^2+n^2$. Opposite parity keeps the triple primitive.
\end{hint}

\section{Solutions}

\begin{solution}{ex:areaper}
\idea Write the condition as two equations, $a^2+b^2=c^2$ and $\frac{ab}2=a+b+c$, and eliminate $c$: substitute $c=\frac{ab}2-a-b$ in the first. After dividing by $ab$ what remains is linear in each of $a,b$, and such equations factor (``Simon's favourite factoring trick''): $ab-4a-4b+8=0$ is $(a-4)(b-4)=8$. Then list the divisors of $8$.
\writeup With legs $a,b$ and hypotenuse $c=\frac{ab}2-a-b$, the equation $c^2=a^2+b^2$ becomes $\frac{a^2b^2}4-a^2b-ab^2+2ab=0$; dividing by $ab$ and multiplying by $4$, $ab-4a-4b+8=0$, that is $(a-4)(b-4)=8$. If both factors are negative, they are among $(-1,-8),(-2,-4),(-4,-2),(-8,-1)$, which give $a\le0$ or $b\le0$. So $\{a-4,b-4\}$ is $\{1,8\}$ or $\{2,4\}$: $(a,b)=(5,12)$ or $(6,8)$, with $c=13$ and $10$. Both work: areas $30$ and $24$, perimeters $30$ and $24$.
\end{solution}

\begin{solution}{ex:x2y23z2}
\idea A homogeneous equation with a non-trivial solution has a primitive one, so look for a prime that forces all variables to be divisible by it. Modulo $3$, $x^2+y^2\equiv0$ forces $x\equiv y\equiv0$ (squares are $0,1$). Then $9\mid3z^2$ gives $3\mid z$: descent.
\writeup Suppose $(x,y,z)\ne(0,0,0)$ is a solution; dividing by $\gcd(x,y,z)$ we may assume it is $1$. Modulo $3$ the squares are $0$ and $1$, and $x^2+y^2\equiv0$ forces $x\equiv y\equiv0\pmod3$. Then $9\mid x^2+y^2=3z^2$, so $3\mid z^2$ and $3\mid z$: $3$ divides all three, a contradiction. A rational point $(\frac XZ,\frac YZ)$ on $x^2+y^2=3$ (common denominator $Z\ge1$) would give $X^2+Y^2=3Z^2$ with $Z\ne0$.
\end{solution}

\begin{solution}{ex:catalan23}
\idea Small cases give $(1,0)$ and $(2,1)$. For $y\ge1$, reduce modulo $3$: $2^x\equiv1$ forces $x$ even, and then $2^x-1=(2^{x/2}-1)(2^{x/2}+1)$ is a difference of squares; two powers of $3$ that differ by $2$ can only be $1$ and $3$.
\writeup If $y=0$, $2^x=2$ and $x=1$. Let $y\ge1$. Modulo $3$, $2^x\equiv1$, and $2^x\equiv(-1)^x$, so $x$ is even, $x=2m$ with $m\ge1$. Then $(2^m-1)(2^m+1)=3^y$, so both factors are powers of $3$ (unique factorisation), and they differ by $2$. The only powers of $3$ differing by $2$ are $1$ and $3$ (if both are $\ge3$, their difference is divisible by $3$). So $2^m=2$, $m=1$, $x=2$, $y=1$. The solutions are $(1,0)$ and $(2,1)$.
\end{solution}

\begin{solution}{ex:mordell5}
\idea Imitate Theorem~\ref{thm:mordell7}: make a sum of two squares appear by adding a square to both sides: $y^2+4=x^3-1=(x-1)(x^2+x+1)$. A prime $q\equiv3\pmod4$ cannot divide $y^2+2^2$ unless it divides both $y$ and $2$. So we need a factor $\equiv3\pmod4$; parity considerations modulo $4$ decide which $x$ can occur.
\writeup Modulo $4$: if $x$ is even, $y^2\equiv-5\equiv3$, impossible. If $x\equiv3\pmod4$, $y^2\equiv27-5\equiv2$, impossible. So $x\equiv1\pmod4$. Then $y^2+4=x^3-1=(x-1)(x^2+x+1)$ and $x^2+x+1\equiv1+1+1=3\pmod4$, with $x^2+x+1\ge3$ (as $x^2+x+1>0$ always, and it is $\equiv3$). So $x^2+x+1$ has a prime factor $q\equiv3\pmod4$, and $q\mid y^2+2^2$. If $q\nmid y$, then $\left(2y^{-1}\right)^2\equiv-1\pmod q$, impossible; so $q\mid y$, then $q\mid4$, impossible for odd $q$.
\end{solution}

\begin{solution}{ex:eisentri}
\idea The conic has the obvious rational point $(-1,0)$, so apply the chord method: intersect with $y=t(x+1)$, factor out $x+1$, and solve the remaining linear equation. For triangles, $c^2=a^2+ab+b^2$ means $(\frac ac,\frac bc)$ is a positive rational point on the conic; write $t=\frac nm$ and clear denominators, as for Pythagorean triples.
\writeup With $y=t(x+1)$:
\[\begin{gathered}x^2-1+tx(x+1)+t^2(x+1)^2=0,\qq\text{i.e.}\quad(x+1)\left(x-1+tx+t^2(x+1)\right)=0.\end{gathered}\]
Besides $(-1,0)$, the points are
\[
\begin{gathered}x=\frac{1-t^2}{1+t+t^2},\qq y=t(x+1)=\frac{2t+t^2}{1+t+t^2},\qq t\in\Q,\end{gathered}
\]
and the vertical line $x=-1$ adds $(-1,1)$. For $a,b,c>0$ with $c^2=a^2+ab+b^2$, the point $(\frac ac,\frac bc)$ has positive coordinates, which requires $0<t<1$; with $t=\frac nm$, $0<n<m$, the triangle is proportional to
\[
\begin{gathered}a=m^2-n^2,\qq b=2mn+n^2,\qq c=m^2+mn+n^2.\end{gathered}
\]
For example $(m,n)=(2,1)$ gives $(3,5,7)$: $9+15+25=49$. Every such triangle is a rational multiple of one of these. To turn this into a list with no omitted integer triangles, divide the three displayed integers by their common gcd $g$, then take every positive integer multiple of the resulting primitive triple. Indeed, if a rational multiple of a primitive integer triple has integer coordinates, its denominator in lowest terms divides all three coordinates and must therefore be $1$.
\end{solution}

\begin{solution}{ex:markovcoprime}
\idea Look for a quantity preserved by the move $x\mapsto3yz-x$: the gcd of the triple is one, since $x'\equiv-x$ modulo any common divisor of $y,z$. As $(1,1,1)$ has gcd $1$, every triple does (Theorem~\ref{thm:markov}). Then a prime dividing two of the numbers divides the square of the third, by the equation. For the list, walk down the tree.
\writeup If $d\mid y$ and $d\mid z$, then $d\mid x$ iff $d\mid x'=3yz-x$; so $\gcd(x,y,z)=\gcd(x',y,z)$, and permutations do not change it either. By Theorem~\ref{thm:markov}, every triple has the gcd of $(1,1,1)$, namely $1$. If a prime $p$ divides $x$ and $y$, then $p^2\mid3xyz-x^2-y^2=z^2$, so $p\mid z$: contradiction. Starting from $(1,1,1)$ and applying moves to the largest-but-new element: $(1,1,2)$, $(1,2,5)$, $(1,5,13)$, $(2,5,29)$, $(1,13,34)$, $(1,34,89)$, $(2,29,169)$, $(5,13,194)$; the next ones exceed $200$ (for instance $(1,89,233)$, $(5,29,433)$). Every triple with maximum $\le200$ appears, because by the proof of Theorem~\ref{thm:markov} each triple other than $(1,1,1)$ has a unique parent with a smaller maximum.
\end{solution}

\begin{solution}{ex:imo88}
\idea For fixed $b$ and $k$, the equation $a^2-kb\,a+b^2-k=0$ is quadratic in $a$, so Vi\`ete gives a second solution $a'=kb-a=\frac{b^2-k}a$ with the same $k$. If $k$ is not a square, this jump should produce a smaller solution, and a minimal counterexample gives a contradiction. The places where the descent could stop are $a'<0$ (impossible by a sign check), $a'=0$ (which makes $k=b^2$, a square), and that is exactly what we want.
\writeup Suppose $k$ is not a square, and among all pairs of non-negative integers $(a,b)$ with $\frac{a^2+b^2}{ab+1}=k$, choose one with $a+b$ minimal, and $a\ge b$. If $b=0$ then $k=a^2$, a square; so $b\ge1$. The number $a'=kb-a=\frac{b^2-k}a$ is an integer and satisfies $a'^2-kb\,a'+b^2-k=0$. If $a'<0$, then $a'^2-kba'+b^2-k\ge1+kb+b^2-k>0$, impossible; so $a'\ge0$. If $a'=0$, then $b^2=k$, a square, excluded. So $a'\ge1$ and $(a',b)$ is again a pair with quotient $k$; and $a'=\frac{b^2-k}a<\frac{b^2}a\le a$. Then $a'+b<a+b$ contradicts minimality. Hence $k$ is a square. (Example: $(a,b)=(8,30)$ gives $k=\frac{964}{241}=4$.)
\end{solution}

\begin{solution}{ex:t1-triple}
\idea Substitute into Theorem~\ref{thm:pyth} and test the largest entry against the sum of the other two squares.
\writeup The pair $(2,1)$ gives $3,4,5$. The pair $(4,1)$ gives $15,8,17$, and $225+64=289$. The pair $(3,2)$ gives $5,12,13$, and $25+144=169$. All three pairs are coprime and of opposite parity.
\end{solution}

\section*{Sources and further reading}

\begin{itemize}
\item L.~J.~Mordell, \emph{Diophantine Equations}, Academic Press, 1969. The classic catalogue of methods, including Mordell equations $y^2=x^3+k$.
\item T.~Andreescu, D.~Andrica and I.~Cucurezeanu, \emph{An Introduction to Diophantine Equations}, Birkh\"auser, 2010. Elementary methods with many solved problems, Vieta jumping included.
\item M.~Aigner, \emph{Markov's Theorem and 100 Years of the Uniqueness Conjecture}, Springer, 2013. Everything about Markov triples and the open conjecture.
\item J.~H.~Silverman and J.~Tate, \emph{Rational Points on Elliptic Curves}, 2nd ed., Springer, 2015, Chapter~1. The chord method on cubics and the group law, the next step after this chapter.
\item Y.~V.~Matiyasevich, \emph{Hilbert's Tenth Problem}, MIT Press, 1993. Why no algorithm can decide every Diophantine equation: a meeting point of number theory and computability.
\item S.~Singh, \emph{Fermat's Last Theorem}, Fourth Estate, 1997. Fermat's marginal note (c.~1637) and Wiles's proof (1994, published in Annals of Math.\ 141, 1995).
\end{itemize}
